% Figura: (a) formación de un sistema de N cargas puntuales trayéndolas una a una
% desde el infinito; (b) ejemplo de las cuatro cargas en un cuadrado.
\begin{figure}[t]
\centering
\begin{subfigure}[b]{0.5\textwidth}\centering
\begin{tikzpicture}[scale=1.0]
  % ejes
  \coordinate (O) at (0,0);
  \draw[tenue, line width=0.45pt, -{Stealth[length=1.6mm]}] (O) -- (2.9,0) node[below, etiqt] {$y$};
  \draw[tenue, line width=0.45pt, -{Stealth[length=1.6mm]}] (O) -- (0,2.4) node[left, etiqt] {$z$};
  \draw[tenue, line width=0.45pt, -{Stealth[length=1.6mm]}] (O) -- (-1.3,-1.0) node[left, etiqt] {$x$};
  % cargas ya colocadas
  \foreach \p/\n in {{(1.3,1.7)}/1,{(2.3,0.9)}/2,{(0.9,-0.9)}/3} {
    \fill[cargapos] \p circle (0.075);
    \draw[campoP, line width=0.6pt, -{Stealth[length=1.6mm]}] (O) -- \p;
  }
  \node[etiq, anchor=south west] at (1.35,1.72) {$q_1$};
  \node[etiq, anchor=west] at (2.38,0.92) {$q_2$};
  \node[etiq, anchor=north west] at (0.95,-0.95) {$q_3$};
  \node[etiq, text=campoP] at (0.42,1.05) {$\vr_1$};
  \node[etiq, text=campoP] at (1.42,0.3) {$\vr_2$};
  \node[etiq, text=campoP] at (0.25,-0.62) {$\vr_3$};
  % carga N trayéndose desde el infinito
  \fill[cargapos] (-1.4,1.55) circle (0.075);
  \node[etiq, anchor=south] at (-1.4,1.62) {$q_N$};
  \draw[campoE, line width=0.6pt, dashed, -{Stealth[length=1.8mm]}] (-2.9,2.25) -- (-1.5,1.6);
  \node[etiqt, anchor=south] at (-2.55,2.3) {desde $\infty$};
  \draw[campoP, line width=0.6pt, -{Stealth[length=1.6mm]}] (O) -- (-1.33,1.48);
  \node[etiq, text=campoP] at (-0.95,0.62) {$\vr_N$};
  \node[etiqt] at (1.9,-0.25) {$\cdots$};
\end{tikzpicture}
\caption{Formación del sistema carga a carga}
\end{subfigure}\hfill
\begin{subfigure}[b]{0.44\textwidth}\centering
\begin{tikzpicture}[scale=1.25]
  \draw[tenue, line width=0.45pt, -{Stealth[length=1.6mm]}] (-0.3,0) -- (2.2,0) node[below, etiqt] {$x$};
  \draw[tenue, line width=0.45pt, -{Stealth[length=1.6mm]}] (0,-0.3) -- (0,2.0) node[left, etiqt] {$y$};
  \draw[equipot] (0,0) -- (1.5,0) -- (1.5,1.5) -- (0,1.5) -- cycle;
  \draw[acentoB, line width=0.55pt] (0,0) -- (1.5,1.5);
  \node[etiqf, text=acentoB, font=\footnotesize, rotate=45] at (0.62,0.82) {$\sqrt2\,a$};
  \foreach \p in {(0,0),(1.5,0),(1.5,1.5),(0,1.5)} \fill[cargapos] \p circle (0.07);
  \node[etiq, anchor=north east] at (-0.03,-0.03) {$q$};
  \node[etiq, anchor=north] at (1.5,-0.08) {$q$};
  \node[etiq, anchor=south west] at (1.53,1.53) {$q$};
  \node[etiq, anchor=east] at (-0.08,1.5) {$q$};
  \draw[cota] (0,-0.55) -- (1.5,-0.55) node[midway, below, etiq] {$a$};
\end{tikzpicture}
\caption{Cuatro cargas en un cuadrado}
\end{subfigure}
\caption{Energía de un sistema de cargas puntuales. (a) El sistema se forma trayendo las cargas una a una desde el
infinito, donde están en reposo e infinitamente alejadas entre sí; la carga $q_N$ debe vencer el campo de las $N-1$ ya
colocadas. (b) Cuatro cargas iguales en los vértices de un cuadrado de lado $a$ (\cref{ej:cuadrado}): cada carga
tiene dos vecinas a distancia $a$ y una a distancia $\sqrt2\,a$.}
\label{fig:cargas}
\end{figure}
